\section{Introduction}
\label{sec:intro}

Remaining useful life (RUL) models are usually judged by their accuracy near the
end of an engine's life, because that is when a maintenance decision is made.
Yet many models produce a sequence of predictions from overlapping windows and
aggregate them with a fixed rule or a learned gate. The aggregation step is often
evaluated on a mixture of early and late cuts, and a single average can hide a
change in ranking across the engine life.

This paper asks a narrow question: \emph{does the benefit of multi-window
aggregation depend on where in the engine life the prediction is made?} We study
the question on the C-MAPSS benchmark \cite{saxena2008} with a frozen backbone and
a fixed evaluation protocol. The protocol separates five cut fractions
(0.40, 0.55, 0.70, 0.85, 0.95), keeps validation and test engines disjoint, and
reports engine-mean and pooled RMSE separately.

Our first result is positive for a learned gate. LTAG aligns the predictions of
several windows to the current cut and learns a gate that mixes the age-aligned
median with the single-window prediction. On the frozen B0 backbone, LTAG lowers
official test RMSE on all four subsets, with 15/15 seed-subset wins and a mean
improvement of 0.391 cycles (Table~\ref{tab:test}).

Our second result is a boundary. On a validation average over all five cuts, two
simple alternatives are stronger: a context-aware residual calibration
(K1-Calib) and a linear stack (LS-ctx). This is not a small effect: K1-Calib and
LS-ctx beat LTAG on all four subsets (Table~\ref{tab:validation}). A naive reading
would reject LTAG. A cut-wise decomposition shows why the two results are
compatible. The calibration methods are much stronger at early cuts, where the
error is large and the validation mean is dominated by those cuts. At 0.85--0.95,
LTAG is competitive with or better than the calibration methods
(Table~\ref{tab:cutwise}, Figure~\ref{fig:cutwise}). The official test uses one
window per engine near end-of-life, so it is closer to the late cuts than to the
validation mean.

Our third result is negative for a measurement-model front-end. An earlier
sensor-reconstruction module (RHSE-TSH) looked effective against a last-value
baseline. With fair backtracked hold and median baselines that use only observed
data, the reconstruction matches or loses on bias and drift and only helps on
stuck sensors (Table~\ref{tab:rhse}). We therefore remove it from the method and
keep the finding as a boundary.

The contributions are:
\begin{enumerate}
\item A cut-controlled evaluation protocol for multi-window RUL aggregation that
reports engine-mean and pooled RMSE separately and never mixes estimators.
\item A positive test result for LTAG on a frozen backbone: 15/15 seed-subset
wins, mean test RMSE improvement 0.391 cycles (Section~\ref{sec:exp}).
\item A cut-wise diagnosis showing that calibration wins early cuts while LTAG
is competitive or better near end-of-life; the validation mean is not a proxy
for the end-of-life test (Section~\ref{sec:cutwise}).
\item A negative result for measurement-model sensor reconstruction under fair
temporal baselines; the module is retained only as a limitation
(Section~\ref{sec:rhse}).
\end{enumerate}

We do not claim a new state of the art. The paper's value is the cut-controlled
evidence and the boundary: a learned gate can be the right choice at end-of-life
even when a simple calibration wins the all-cut average.

The rest of the paper describes the protocol (Section~\ref{sec:setup}), the LTAG
method (Section~\ref{sec:method}), the test and validation results
(Section~\ref{sec:exp}), the cut-wise analysis (Section~\ref{sec:cutwise}), the
negative sensor-reconstruction result (Section~\ref{sec:rhse}), and the
limitations (Section~\ref{sec:conclusion}).
